\documentclass[a4paper,12pt]{article}

% Packages
\usepackage[utf8]{inputenc}
\usepackage[T1]{fontenc}
\usepackage{listings}
\usepackage{amsmath}
\usepackage{graphicx}
\usepackage{xcolor}
\usepackage{geometry}
\usepackage{titlesec}
\usepackage{booktabs}
\usepackage{enumitem}
\usepackage{fancyhdr}
\usepackage{tcolorbox}
\usepackage{lmodern}
\usepackage{array} 
\usepackage{caption}



% Page setup
\geometry{a4paper, margin=2.5cm, marginpar=1.5cm}
\setlength{\parindent}{0pt}
\setlength{\parskip}{0.8em}
\renewcommand{\baselinestretch}{1.1}

% Custom colors
\definecolor{titleblue}{RGB}{0,76,151}
\definecolor{sectiongray}{RGB}{80,80,80}
\definecolor{codebackground}{RGB}{245,245,245}

% Header/footer
\pagestyle{fancy}
\fancyhf{}
\renewcommand{\headrulewidth}{0.4pt}
\renewcommand{\footrulewidth}{0.4pt}
\fancyhead[R]{\textcolor{titleblue}{Microprocessors - Lab 1 Report}}
\fancyfoot[C]{\textcolor{titleblue}{\thepage}}
\setlength{\headheight}{14.5pt}
\addtolength{\topmargin}{-2.5pt}

% Title formatting
\titleformat{\section}{\Large\bfseries\sffamily\color{titleblue}}{\thesection}{1em}{}
\titleformat{\subsection}{\large\bfseries\sffamily\color{titleblue}}{\thesubsection}{1em}{}

% Custom section boxes
\newenvironment{sectionbox}[1]{%
    \tcolorbox[
        colback=titleblue!5,
        colframe=titleblue!50,
        arc=3mm,
        title=#1,
        fonttitle=\bfseries\sffamily,
        before skip=1.2em,
        after skip=1.2em
    ]
}{\endtcolorbox}

% Code listing style
\lstdefinestyle{mystyle}{
    backgroundcolor=\color{codebackground},   
    commentstyle=\color{green!50!black},
    keywordstyle=\color{titleblue},
    numberstyle=\tiny\color{sectiongray},
    stringstyle=\color{purple!70!black},
    basicstyle=\ttfamily\footnotesize,
    breakatwhitespace=false,         
    breaklines=true,                 
    captionpos=b,                    
    keepspaces=true,                 
    numbers=left,                    
    numbersep=8pt,                  
    showspaces=false,                
    showstringspaces=false,
    showtabs=false,                  
    tabsize=2,
    frame=single,
    rulecolor=\color{titleblue!30},
    frameround=tttt
}

\lstset{style=mystyle}

% Document Information
\title{
    \vspace*{2cm}
    \Huge\textbf{\textsf{Microprocessors Laboratory Report}}\\
    \vspace{0.4cm}
    \Large\textsf{Lab 1: Assembly Programming \& UART Communication}
    \vspace{1cm}
    \rule{\textwidth}{1pt}
}

\author{
    \Large\begin{tabular}{rl}
        \textbf{Fraidakis Ioannis}\\
    \end{tabular}
}

\date{\textsf{April 2025}}


\begin{document}
\maketitle

\vspace{0.5cm}

\begin{sectionbox}{Executive Summary}
This laboratory exercise explores low-level microcontroller programming through three ARM Cortex-M4 assembly language implementations, seamlessly integrated with C for system-level operation. The project focuses on:

\begin{itemize}[leftmargin=2em, itemsep=0.5em]
    \item \textbf{Character-based Hash Algorithm:} Implementation of a custom string hashing function with character-specific processing logic and then optimize it to achieve a 43.5\% performance improvement
    
    \item \textbf{Recursive Computation:} Implementation of a recursive Fibonacci function
        
    \item \textbf{Hardware Communication:} Integration with UART peripheral for real-time data exchange between the microcontroller and external systems
\end{itemize}

The implementations demonstrate critical concepts in embedded systems development, including mixed-language programming, hardware peripheral control, and algorithm optimization techniques specifically designed for resource-constrained environments. Each solution addresses different aspects of microcontroller programming, from computational efficiency to memory management and I/O operations.
\end{sectionbox}





\begin{sectionbox}{Project Structure}
The project is organized into a modular structure with dedicated folders for each problem solution and multiple input modes:

\subsection*{Main Entry Points}
Multiple C files provide different input methods for testing:
\begin{itemize}[leftmargin=2em, itemsep=0.5em]
    \item \textbf{hardcoded\_input.c} - Input is provided directly in the code
    \item \textbf{UART\_input.c} - Receives input via UART hardware interface
    \item \textbf{scanf\_input.c} - Uses Keil's serial console for input
    \item \textbf{test\_optimization.c} - Performance benchmark harness for Question 1
    \item \textbf{return\_ptr.c} - Used to run \texttt{fib\_return\_ptr.s} variant for Question 3
\end{itemize}


\begin{tcolorbox}[colback=titleblue!3,colframe=titleblue!30,arc=2mm]
\textbf{Execution Note:} The project is configured to use hardcoded\_input.c by default. Other entry points can be enabled by uncommenting the desired file and commenting out the currently active one.
\end{tcolorbox}


\subsection*{Implementation Folders}
\begin{itemize}[leftmargin=2em, itemsep=0.5em]
    \item \textbf{Question\_1/} - String hashing implementations:
    \begin{itemize}
        \item \texttt{hash.s} - Basic implementation
        \item \texttt{hash\_optimized.s} - Optimized version with performance improvements
        \item \texttt{min\_registers.s} - Uses only scratch registers (no stack usage)
    \end{itemize}

    \item \textbf{Question\_2/} - Digit reduction algorithm:
    \begin{itemize}
        \item \texttt{sum\_mod.s} - Straightforward implementation
    \end{itemize}

    \item \textbf{Question\_3/} - Fibonacci implementations:
    \begin{itemize}
        \item \texttt{fibonacci.s} - Standard recursive implementation
        \item \texttt{fib\_return\_ptr.s} - Returns address of result, stored in a local variable
    \end{itemize}

    \item \textbf{Bonus/} - Contains the solution for the bonus question
\end{itemize}
\end{sectionbox}

\begin{sectionbox}{Technical Implementation}
\subsection*{Core Function Analysis}

\begin{itemize}[leftmargin=2em]

    \item \textbf{String Hasher (hash\_final.s)} - Character processing logic:
    \begin{lstlisting}[language={[x86masm]Assembler}]
    ; Digit conversion
    LDR R8, =digit_table  ; Load lookup table
    LDRB R7, [R8, R7]     ; Get table value
    
    ; Uppercase handling
    ADDLE R5, R5, R6, LSL #1  ; 2*ASCII via barrel shifter
    
    ; Lowercase calculation
    SUB R7, R6, #'a'       ; Offset from 'a'
    MLA R5, R7, R7, R5     ; (offset)^2 + hash
    \end{lstlisting}

    \textbf{Processing Logic:}
    \begin{itemize}
        \item \textit{Digits:} Mapped to predefined values (e.g., '0'→5, '1'→12) via \texttt{digit\_table}
        \item \textit{Uppercase:} Add 2×ASCII value via barrel shifter (LSL \#1)
        \item \textit{Lowercase:} Add (char-'a')² value using MLA instruction (fusing multiply and add)
        \item \textit{Special Chars:} Skipped in processing
    \end{itemize}

\item \textbf{Digit Reducer (\texttt{sum\_mod.s})} – Core reduction logic:
\begin{lstlisting}[language={[x86masm]Assembler}]
    ; Digit summation
    MOV R5, #10        ; Set divisor for decimal base
    UDIV R6, R4, R5    ; R6 = R4 / 10
    MLS R7, R5, R6, R4 ; R7 = R4-(R5*R6) -> extract last digit
    
    ; Modulo 7 operation
    MOV R5, #7         ; Set divisor for modulo
    UDIV R6, R4, R5    ; R6 = R4 / 7
    MLS R7, R5, R6, R4 ; R7 = R4 - (R5 * R6) -> R4 mod 7
\end{lstlisting}

\textbf{Algorithm:}
    \begin{itemize}
    \item Extracts each digit of the hash value (e.g., 1234 → 1, 2, 3, 4) by repeatedly dividing by 10 and accumulating the sum (1234 → 1+2+3+4=10).
    \item Computes modulo 7 of the sum (e.g., 10 \% 7 = 3). 
    \end{itemize}
\end{itemize}
\end{sectionbox}









\begin{sectionbox}{Technical Implementation (continued)}
\subsection*{Core Function Analysis}

\begin{itemize}[leftmargin=2em]


    \item \textbf{Fibonacci (fibonacci.s)} - Recursive implementation with stack management:
    \begin{lstlisting}[language={[x86masm]Assembler}]
    ; Base case handling
    CMP R0, #0          ; Check n=0
    BXEQ LR             ; Return 0 (R0 is already 0)
    CMP R0, #1          ; Check n=1
    BXEQ LR             ; Return 1 (R0 is already 1)
    
    ; Recursive decomposition
    PUSH {LR, R0}       ; Save return address and n
    SUB R0, R0, #1      ; n-1
    BL fibonacci        ; fib(n-1)
    POP {R1}            ; Restore n
    PUSH {R0}           ; Save fib(n-1)
    SUB R0, R1, #2      ; n-2
    BL fibonacci        ; fib(n-2)
    POP {R1}            ; R1 = fib(n-1)
    ADD R0, R0, R1      ; R0 = fib(n-1) + fib(n-2)
    POP {PC}            ; Return
    \end{lstlisting}

    \textbf{Mechanism:} Implements the recursive definition fib(n) = fib(n-1) + fib(n-2):
    \begin{itemize}
        \item Preserves return address (LR) and registers using PUSH/POP
        \item Uses R0 for both input parameter and return value
        \item Uses conditional return for base cases to minimize branching
    \end{itemize}
\end{itemize}
\end{sectionbox}


\begin{sectionbox}{Testing \& Validation}
    \subsection*{Testing Methodology}
    To validate our implementations, we performed comprehensive testing using both static test vectors and dynamic inputs. For performance evaluation, we used the Keil $\mu$Vision simulator with "Performance analyzer" enabled to measure execution time of the optimized versus the baseline hash implementation.
    
\subsection*{Validation Results}
\begin{center}
\begin{tabular}{@{}lccccc@{}}
\toprule
\textbf{Input} & \textbf{Raw Hash} & \textbf{Reduced} & \textbf{Fib} & \textbf{Checksum} & \textit{Characteristics} \\
\midrule
"" (empty)        & 0    & 0 & 0  & 0  & Edge case \\
"Keil uVision"     & 1752  & 1 & 1  & 90  & Alphabetic \\
"STM32-F411RE"    & 983 & 6 & 8 & 3  & Mixed case \\
"Arm 2025"       & 601  & 0 & 0  & 123  & Alphanumeric \\

\bottomrule
\end{tabular}
\end{center}
\end{sectionbox}





\begin{sectionbox}{Optimization Techniques}

\subsection*{Optimization Strategies}

Key optimizations implemented in the hash function:

\begin{itemize}[leftmargin=2em]
    
    \item \textbf{Single-Pass Processing:} Integrated string length calculation during character processing (eliminated separate traversal)
    
    \item \textbf{Conditional execution:} Implementing uppercase character processing with conditional execution (LE suffix) to avoid branch penalties:
    \begin{lstlisting}[language={[x86masm]Assembler}]
    CMP     R6, #'Z'           ; Compare with 'Z'
    ADDLE   R5, R5, R6, LSL #1 ; Conditionally execute if <='Z'
    BLE     loop_start
    \end{lstlisting}
    
    \item \textbf{Barrel shifter utilization:} Employing the ARM barrel shifter for efficient multiplication by powers of 2:
    \begin{lstlisting}[language={[x86masm]Assembler}]
    ADDLE R5, R5, R6, LSL #1 ; Instead of separate shift
    \end{lstlisting}
    
    \item \textbf{Combined operations:} Using MLA (Multiply-Accumulate) to perform multiplication and addition in a single instruction:
    \begin{lstlisting}[language={[x86masm]Assembler}]
    MLA R5, R6, R6, R5  ; R5 = (R6*R6) + R5 in one cycle
    \end{lstlisting}
    
    \item \textbf{Efficient character classification:} Implementing a single-branch lowercase character check:
    \begin{lstlisting}[language={[x86masm]Assembler}]
    SUB r0, r0, #'a'    ; Offset from 'a'
    CMP r0, #25         ; Check if in range 0-25 ('a' to 'z')
    BHI not_lowercase   ; Branch if higher (unsigned)
    \end{lstlisting}    

    \item \textbf{Post-index LDRB}: Used to combine load and increment in one instruction:
    \begin{lstlisting}[language={[x86masm]Assembler}]
    LDRB R6, [R4], #1  ; Instead of separate ADD for R4
    \end{lstlisting}


\end{itemize}

\begin{tcolorbox}[colback=titleblue!3,colframe=titleblue!30,arc=2mm]
\textbf{Performance Optimizations:} Our optimized hash implementation achieved a 
\textbf{43.5\% speedup} over the baseline implementation. The benchmark was performed on 
a 1000-character string with an equal distribution of digits, uppercase and lowercase letters, and some symbols.
\end{tcolorbox}
\end{sectionbox}






\begin{sectionbox}{Conclusion \& Learning Outcomes}
This lab provided valuable experience in several areas of embedded programming:

\begin{itemize}[leftmargin=2em]
    \item \textbf{ARM Assembly Mastery:} Gained practical experience with ARM instruction set architecture, including efficient use of specialized instructions.
    
    \item \textbf{Register Optimization:} Understanding of register allocation constraints, a critical skill for interrupt handlers and resource-constrained environments
    
    \item \textbf{Optimization Techniques:} Learned to identify and implement performance optimizations for resource-constrained environments, achieving significant speed improvements through instruction selection and branch elimination
    
    \item \textbf{Mixed-Language Programming:} Developed skill in integrating C and assembly code, including proper parameter passing and return value handling    
\end{itemize}

% The techniques learned in this laboratory exercise provide a foundation for more complex embedded systems development, particularly in scenarios where performance optimization, register constraints, and hardware interaction are critical requirements.
\end{sectionbox}

\begin{sectionbox}{Future Work}
Based on our experience and performance analysis, several optimization opportunities remain for future exploration:

\begin{itemize}[leftmargin=2em]
    \item \textbf{Elimination of Division Instructions:}
    Replace \texttt{UDIV} operations with multiplication-based division for constant divisors (e.g. 10). For example:
    \begin{lstlisting}[language={[x86masm]Assembler}]
    ; Instead of:
    MOV R5, #10
    UDIV R6, R4, R5
    
    ; Division by 10 via reciprocal
    LDR R5, =0xCCCCCCCD ; Magic constant for 1/10 in Q32 format
    UMULL R6, R7, R4, R5 ; R7 = (R4 * magic) >> 32 = R4 / 10
    MOV R7, R7, LSR #3   ; Complete the division
    \end{lstlisting}

    \begin{tcolorbox}[colback=titleblue!3,colframe=titleblue!30,arc=2mm]
    \textbf{Impact on Cortex-M4:}  
    While the M4 implements \texttt{UDIV} in hardware (2-12 cycles, divisor-dependent), reciprocal multiplication reduces worst-case execution time (WCET). 
    \end{tcolorbox}
        
    \item \textbf{Loop unrolling:} Apply strategic loop unrolling in the string processing loop to reduce branch prediction overhead for long input strings
    \item \textbf{Register allocation refinement:} Further tune register usage patterns for improved locality and reduced memory operations
    \item \textbf{Hybrid Digital Root:}  
    Combine digit summation with modulo 9 properties (e.g., \texttt{digital\_root = 1 + (num - 1) \% 9}) for single-step digit addition

\end{itemize}
\end{sectionbox}


\end{document}
